Throughout this survey, we have systematically explored the intricate landscape of factuality issues within large language models (LLMs). We began by defining the concept of factuality (Sec \ref{sec:definition}) and proceeded to discuss its broader implications (Sec \ref{sec:impact}). Our journey took us through the multifaceted realm of factuality evaluation, encompassing benchmarks (Sec \ref{sec:eval_benchmark}), metrics (Sec \ref{sec:eval_metrics}), specific evaluation studies (Sec \ref{sec:factuality-eval}), and domain-specific evaluations (Sec \ref{sec:domain-eval}). We then delved deeper, probing the intrinsic mechanisms that underpin factuality in LLMs (Sec \ref{sec:analysis}). Our exploration culminated in the discussion of enhancement techniques, both for standalone LLMs (Sec \ref{sec:enhace_pure}) and retrieval-augmented LLMs (Sec \ref{sec:enhace_rag}), with a special focus on domain-specific LLM enhancements (Sec \ref{sec:domain_llms}).

Despite the advancements detailed in this survey, several challenges loom large. The evaluation of factuality remains an intricate puzzle, complicated by the inherent variability and nuances of natural languages. The core processes governing how LLMs store, update, and produce facts are yet not fully revealed. And while certain techniques, like continual training and retrieval, show promise, they are not without limitations.
Looking ahead, the quest for fully factual LLMs presents both challenges and opportunities. Future research might delve deeper into understanding the neural architectures of LLMs, develop more robust evaluation metrics, and innovate on enhancement techniques. As LLMs become increasingly integrated into our digital ecosystem, ensuring their factual reliability will remain paramount, with implications that impact across the AI community and beyond.